\clearpage\pagestyle{otMain}\setlength{\otBlocEspacement}{2.5mm}\fontsize{11}{15}\selectfont\corrigesChapitre{8}{Lettres, formules et motifs}
\begin{blocLivre}{cor-8-1}\begin{corrige}{1}{Double et successeur}\demarche{Traduire les mots en opérations.}
Le double est $\res{2x}$ ; le successeur est $\res{x+1}$. Les opérations sont différentes.
\end{corrige}\end{blocLivre}
\begin{blocLivre}{cor-8-2}\begin{corrige}{2}{Lire un produit}\demarche{Utiliser les conventions.}
On écrit $\res{5a}$ et $\res{b^2}$. On ne remplace pas $b\times b$ par $2b$, qui est une somme.
\end{corrige}\end{blocLivre}
\begin{blocLivre}{cor-8-3}\begin{corrige}{3}{Calculer une formule}\demarche{Substituer puis appliquer les priorités.}
$4\times2+3=8+3=\res{11}$.
\end{corrige}\end{blocLivre}
\begin{blocLivre}{cor-8-4}\begin{corrige}{4}{Une aire}\demarche{Relier la formule à la grandeur.}
La formule est $A=5x$. Pour $x=7$, $A=5\times7=\res{35\ \mathrm{cm}^2}$.
\label{fig-cor-8-4}\figureGeo[corrige]{lit-4}
\end{corrige}\end{blocLivre}
\begin{blocLivre}{cor-8-5}\begin{corrige}{5}{Tester une égalité}\demarche{Calculer les deux membres.}
Le membre gauche vaut $3\times4+2=14$, égal au membre droit. L’égalité est $\res{\text{vraie}}$ pour cette valeur.
\end{corrige}\end{blocLivre}
\begin{blocLivre}{cor-8-6}\begin{corrige}{6}{Une valeur décimale}\demarche{Calculer le carré avant le coefficient.}
$1{,}5^2=2{,}25$, donc $2\times2{,}25+1=\res{5{,}5}$.
\end{corrige}\end{blocLivre}
\begin{blocLivre}{cor-8-7}\begin{corrige}{7}{Une somme ou un produit}\demarche{Lire la structure principale.}
$3x+5$ est une somme, valant $\res{11}$. $3(x+5)$ est un produit, valant $\res{21}$. On nomme la dernière opération.
\end{corrige}\end{blocLivre}
\begin{blocLivre}{cor-8-8}\begin{corrige}{8}{Motif d’allumettes}\demarche{Identifier ce qui est commun entre deux motifs.}
Avec $3n+1$, on trouve $\res{4,\ 7,\ 16,\ 61}$. Chaque nouveau carré ajoute trois allumettes, car un côté est partagé.
\label{fig-cor-8-8}\figureGeo[corrige]{lit-8}
\end{corrige}\end{blocLivre}
\begin{blocLivre}{cor-8-9}\begin{corrige}{9}{Un programme verbal}\demarche{Respecter l’ordre dans la formule.}
La somme est faite d’abord : $\res{3(x+4)}$. Pour $x=6$, le résultat vaut $3\times10=\res{30}$.
\end{corrige}\end{blocLivre}
\begin{blocLivre}{cor-8-10}\begin{corrige}{10}{Une formule de tableur}\demarche{Relier une cellule à la valeur d’entrée.}
La formule calcule $4x+3$. Pour 0, 2 et 5, on obtient $\res{3,\ 11,\ 23}$. Après recopie, la référence A2 devient A3 puis A4.
\end{corrige}\end{blocLivre}
\begin{blocLivre}{cor-8-11}\begin{corrige}{11}{Une bordure de points}\demarche{Corriger les éléments comptés deux fois.}
Compter $4n$ compte chaque sommet deux fois. On retire les quatre doublons : $\res{4n-4}$. Pour dix points par côté : $\res{36}$ points.
\label{fig-cor-8-11}\figureGeo[corrige]{lit-11}
\end{corrige}\end{blocLivre}
\begin{blocLivre}{cor-8-12}\begin{corrige}{12}{Un tarif}\demarche{Séparer partie fixe et partie variable.}
Le prix en euros est $\res{8+3h}$. Pour $h=7$, on paie $8+21=\res{29}$ euros. Le coût fixe n’est ajouté qu’une fois.
\end{corrige}\end{blocLivre}
\begin{blocLivre}{cor-8-13}\begin{corrige}{13}{Deux calculs à tester}\demarche{Distinguer test et preuve.}
Pour 2, les deux valent 8. Cela ne suffit pas : pour 1, ils valent 5 et 4. L’égalité n’est donc $\res{\text{pas toujours vraie}}$.
\end{corrige}\end{blocLivre}
\begin{blocLivre}{cor-8-14}\begin{corrige}{14}{Formule d’une suite}\demarche{Tenir compte du numéro de départ.}
On part de 5 puis on ajoute quatre jetons $n-1$ fois : $5+4(n-1)=\res{4n+1}$. L’étape 25 contient $\res{101}$ jetons.
\end{corrige}\end{blocLivre}
\begin{blocLivre}{cor-8-15}\begin{corrige}{15}{Vérifier une vraisemblance}\demarche{Tester plusieurs valeurs puis revenir à la définition.}
Le rectangle mesure 6 cm sur 3 cm : son périmètre est 18 cm. La formule proposée donne aussi 18 ici, donc ce test ne révèle pas l’erreur. Pour $x=1$, le périmètre vaut 12 cm mais la formule donne 9. La bonne formule est $\res{P=2(x+2)+6=2x+10}$.
\label{fig-cor-8-15}\figureGeo[corrige]{lit-15}
\end{corrige}\end{blocLivre}
